% GENERATED FILE. Edit canonical lessons, then run npm run generate:books.
% canonical-source-hash: fnv1a64:7973a9da56465b15
% canonical-lessons: UR-C385-vazifa, UR-C385-farq, UR-C385-mauzu, UR-C385-dakhila, UR-C385-tarjih-dena, UR-R385-review-379-385, UR-R385-review-final-385

\chapter{Choosing What to Study}
\label{ch:ur-choosing-what-to-study}
\hlchaptermodality{385}

\begin{chapteropening}
\textbf{By the end of this chapter:} \emph{I can say what I will study if I get the chance.}

Complete the last lesson of chapter 385: I can say what I will study if I get the chance.
\end{chapteropening}

\section[\texorpdfstring{\ur{وظیفہ} (vazīfā)}{وظیفہ (vazīfā)}]{\ur{وظیفہ} (vazīfā) — a scholarship; a stipend}

\label{lesson:UR-C385-vazifa}



\begin{quote}
Before the new one: say the Urdu for pampering, then the Urdu for to learn by rote.
\end{quote}

\subsection*{You'll want to know: \ur{وظیفہ}}
\textbf{\ur{وظیفہ}} — \emph{vazīfā} — \textquotedblleft{}a scholarship; a stipend\textquotedblright{}.

\begin{quote}
\ur{اگر} \ur{مجھے} \ur{وظیفہ} \ur{ملا} \ur{تو} \ur{میں} \ur{یونیورسٹی} \ur{جاؤں} \ur{گا۔}
\end{quote}

\begin{quote}
\emph{agar mujhe vazīfā milā to maiṅ yūnīvarsiṭī jāūṅ gā.}
\end{quote}

\begin{quote}
If I get a scholarship, I will go to university.
\end{quote}

Masculine.

\begin{cousinweb}
\textbf{\ur{وظیفہ}} (\emph{vazīfā}) is usually traced, through Persian, to Arabic \emph{waẓīfa}, \textquotedblleft{}a duty; an allowance\textquotedblright{}.
\end{cousinweb}

\begin{grammarlens}[title={What you've built}]
One more word for saying what happens if.
\end{grammarlens}

\subsection*{Your turn}
\begin{itemize}
\raggedright
  \item \emph{Say it:} \emph{vazīfā}
  \item \emph{Say it:} \emph{vazīfā}, once more
  \item \emph{Say it:} \emph{agar mujhe vazīfā milā to maiṅ yūnīvarsiṭī jāūṅ gā.} — If I get a scholarship, I will go to university.
  \item \emph{Recall:} say the Urdu for pampering, then the Urdu for to learn by rote, then say \emph{vazīfā} again
\end{itemize}

\begin{tcolorbox}[breakable,colback=teal!4,colframe=teal!35!black,title={Before you move on}]
What does \ur{وظیفہ} mean? (\textquotedblleft{}A scholarship; a stipend\textquotedblright{}.) Say it once more.
\end{tcolorbox}
\section[\texorpdfstring{\ur{فرق} (farq)}{فرق (farq)}]{\ur{فرق} (farq) — a difference}

\label{lesson:UR-C385-farq}



\begin{quote}
Before the new one: say the Urdu for to learn by rote, then the Urdu for a scholarship.
\end{quote}

\subsection*{You'll want to know: \ur{فرق}}
\textbf{\ur{فرق}} — \emph{farq} — \textquotedblleft{}a difference\textquotedblright{}.

\begin{quote}
\ur{اگر} \ur{دونوں} \ur{میں} \ur{فرق} \ur{نہیں} \ur{تو} \ur{سستا} \ur{خریدیے۔}
\end{quote}

\begin{quote}
\emph{agar donoṅ meṅ farq nahīṅ to sastā kharīdiye.}
\end{quote}

\begin{quote}
If there is no difference between them, buy the cheaper one.
\end{quote}

Masculine.

\begin{cousinweb}
\textbf{\ur{فرق}} (\emph{farq}) is usually traced, through Persian, to Arabic \emph{farq}, \textquotedblleft{}separating, a parting\textquotedblright{}.
\end{cousinweb}

\begin{grammarlens}[title={What you've built}]
One more word for saying what happens if.
\end{grammarlens}

\subsection*{Your turn}
\begin{itemize}
\raggedright
  \item \emph{Say it:} \emph{farq}
  \item \emph{Say it:} \emph{farq}, once more
  \item \emph{Say it:} \emph{agar donoṅ meṅ farq nahīṅ to sastā kharīdiye.} — If there is no difference between them, buy the cheaper one.
  \item \emph{Recall:} say the Urdu for to learn by rote, then the Urdu for a scholarship, then say \emph{farq} again
\end{itemize}

\begin{tcolorbox}[breakable,colback=teal!4,colframe=teal!35!black,title={Before you move on}]
What does \ur{فرق} mean? (\textquotedblleft{}A difference\textquotedblright{}.) Say it once more.
\end{tcolorbox}
\section[\texorpdfstring{\ur{موضوع} (mauzū)}{موضوع (mauzū)}]{\ur{موضوع} (mauzū) — a topic; a subject (of study)}

\label{lesson:UR-C385-mauzu}



\begin{quote}
Before the new one: say the Urdu for a scholarship, then the Urdu for a difference.
\end{quote}

\subsection*{You'll want to know: \ur{موضوع}}
\textbf{\ur{موضوع}} — \emph{mauzū} — \textquotedblleft{}a topic; a subject (of study)\textquotedblright{}.

\begin{quote}
\ur{اگر} \ur{یہ} \ur{موضوع} \ur{مشکل} \ur{ہے} \ur{تو} \ur{دوسرا} \ur{چنیے۔}
\end{quote}

\begin{quote}
\emph{agar yih mauzū mushkil hai to dūsrā chuniye.}
\end{quote}

\begin{quote}
If this topic is hard, choose another.
\end{quote}

Masculine.

\begin{cousinweb}
\textbf{\ur{موضوع}} (\emph{mauzū}) is usually traced, through Persian, to Arabic \emph{mawḍūʿ}, \textquotedblleft{}something laid down\textquotedblright{}: the matter set before you.
\end{cousinweb}

\begin{grammarlens}[title={What you've built}]
One more word for saying what happens if.
\end{grammarlens}

\subsection*{Your turn}
\begin{itemize}
\raggedright
  \item \emph{Say it:} \emph{mauzū}
  \item \emph{Say it:} \emph{mauzū}, once more
  \item \emph{Say it:} \emph{agar yih mauzū mushkil hai to dūsrā chuniye.} — If this topic is hard, choose another.
  \item \emph{Recall:} say the Urdu for a scholarship, then the Urdu for a difference, then say \emph{mauzū} again
\end{itemize}

\begin{tcolorbox}[breakable,colback=teal!4,colframe=teal!35!black,title={Before you move on}]
What does \ur{موضوع} mean? (\textquotedblleft{}A topic; a subject (of study)\textquotedblright{}.) Say it once more.
\end{tcolorbox}
\section[\texorpdfstring{\ur{داخلہ} (dākhilā)}{داخلہ (dākhilā)}]{\ur{داخلہ} (dākhilā) — admission (to a school or college)}

\label{lesson:UR-C385-dakhila}



\begin{quote}
Before the new one: say the Urdu for a difference, then the Urdu for a topic.
\end{quote}

\subsection*{You'll want to know: \ur{داخلہ}}
\textbf{\ur{داخلہ}} — \emph{dākhilā} — \textquotedblleft{}admission (to a school or college)\textquotedblright{}.

\begin{quote}
\ur{اگر} \ur{نمبر} \ur{اچھے} \ur{ہوئے} \ur{تو} \ur{کالج} \ur{میں} \ur{داخلہ} \ur{ملے} \ur{گا۔}
\end{quote}

\begin{quote}
\emph{agar nambar achchhe hue to kālij meṅ dākhilā mile gā.}
\end{quote}

\begin{quote}
If the marks are good, there will be a place at college.
\end{quote}

Masculine.

\begin{cousinweb}
\textbf{\ur{داخلہ}} (\emph{dākhilā}) is usually traced, through Persian, to Arabic \emph{dākhil}, \textquotedblleft{}entering, inside\textquotedblright{}: a way in.
\end{cousinweb}

\begin{grammarlens}[title={What you've built}]
One more word for saying what happens if.
\end{grammarlens}

\subsection*{Your turn}
\begin{itemize}
\raggedright
  \item \emph{Say it:} \emph{dākhilā}
  \item \emph{Say it:} \emph{dākhilā}, once more
  \item \emph{Say it:} \emph{agar nambar achchhe hue to kālij meṅ dākhilā mile gā.} — If the marks are good, there will be a place at college.
  \item \emph{Recall:} say the Urdu for a difference, then the Urdu for a topic, then say \emph{dākhilā} again
\end{itemize}

\begin{tcolorbox}[breakable,colback=teal!4,colframe=teal!35!black,title={Before you move on}]
What does \ur{داخلہ} mean? (\textquotedblleft{}Admission (to a school or college)\textquotedblright{}.) Say it once more.
\end{tcolorbox}
\section[\texorpdfstring{\ur{ترجیح} \ur{دینا} (tarjīh denā)}{ترجیح دینا (tarjīh denā)}]{\ur{ترجیح} \ur{دینا} (tarjīh denā) — to prefer; to give priority to}

\label{lesson:UR-C385-tarjih-dena}



\begin{quote}
Before the new one: say the Urdu for a topic, then the Urdu for admission.
\end{quote}

\subsection*{You'll want to know: \ur{ترجیح} \ur{دینا}}
\textbf{\ur{ترجیح} \ur{دینا}} — \emph{tarjīh denā} — \textquotedblleft{}to prefer; to give priority to\textquotedblright{}.

\begin{quote}
\ur{اگر} \ur{انتخاب} \ur{میرا} \ur{ہوتا} \ur{تو} \ur{میں} \ur{تاریخ} \ur{کو} \ur{ترجیح} \ur{دیتا۔}
\end{quote}

\begin{quote}
\emph{agar intikhāb merā hotā to maiṅ tārīkh ko tarjīh detā.}
\end{quote}

\begin{quote}
If the choice were mine, I would pick history.
\end{quote}

Past, with \ur{نے} (he or she …): \ur{اس} \ur{نے} \ur{ترجیح} \ur{دی}, \emph{us ne tarjīh dī}.

\begin{cousinweb}
\textbf{\ur{ترجیح} \ur{دینا}} (\emph{tarjīh denā}) is built from \emph{tarjīh}, usually traced, through Persian, to Arabic \emph{tarjīḥ}, \textquotedblleft{}giving more weight (to one side)\textquotedblright{}, and \textbf{\ur{دینا}}, \textquotedblleft{}to give\textquotedblright{}.
\end{cousinweb}

\begin{grammarlens}[title={What you've built}]
One more word for saying what happens if.
\end{grammarlens}

\subsection*{Your turn}
\begin{itemize}
\raggedright
  \item \emph{Say it:} \emph{tarjīh denā}
  \item \emph{Say it:} \emph{tarjīh denā}, once more
  \item \emph{Say it:} \emph{agar intikhāb merā hotā to maiṅ tārīkh ko tarjīh detā.} — If the choice were mine, I would pick history.
  \item \emph{Recall:} say the Urdu for a topic, then the Urdu for admission, then say \emph{tarjīh denā} again
\end{itemize}

\begin{tcolorbox}[breakable,colback=teal!4,colframe=teal!35!black,title={Before you move on}]
What does \ur{ترجیح} \ur{دینا} mean? (\textquotedblleft{}To prefer; to give priority to\textquotedblright{}.) Say it once more.
\end{tcolorbox}
\section[(dialogue)]{Looking back: from Effort to Choosing What to Study}

\label{lesson:UR-R385-review-379-385}



\begin{quote}
Say the Urdu for hard work and enthusiasm.
\end{quote}

\subsection*{Your turn}
Say these aloud:

\begin{itemize}
\raggedright
  \item hard work, enthusiasm, dedication, without a break, to go over again
  \item a method, practice, nerve, to occur to, to get down to
  \item successful, a feat, pride, a marvel, well done
  \item a prize, an award, a medal, deserving, to tidy up
  \item approval, an arrangement, absolute, to rake in, to fail an exam
  \item upbringing, generous, manners, pampering, to learn by rote
  \item a scholarship, a difference, a topic, admission, to prefer
\end{itemize}

\begin{tcolorbox}[breakable,colback=teal!4,colframe=teal!35!black,title={Before you move on}]
Hard work? (\textbf{\ur{محنت}}, \emph{mehnat}.) Enthusiasm? (\textbf{\ur{جوش}}, \emph{josh}.) Dedication? (\textbf{\ur{لگن}}, \emph{lagan}.) Without a break? (\textbf{\ur{مسلسل}}, \emph{musalsal}.) To go over again? (\textbf{\ur{دہرانا}}, \emph{dohrānā}.) A method? (\textbf{\ur{طریقہ}}, \emph{tarīqā}.) Practice? (\textbf{\ur{مشق}}, \emph{mashq}.) Nerve? (\textbf{\ur{ہمت}}, \emph{himmat}.) To occur to? (\textbf{\ur{سوجھنا}}, \emph{sūjhnā}.) To get down to? (\textbf{\ur{جٹنا}}, \emph{juṭnā}.) Successful? (\textbf{\ur{کامیاب}}, \emph{kāmyāb}.) A feat? (\textbf{\ur{کارنامہ}}, \emph{kārnāmā}.) Pride? (\textbf{\ur{فخر}}, \emph{fakhr}.) A marvel? (\textbf{\ur{کمال}}, \emph{kamāl}.) Well done? (\textbf{\ur{شاباش}}, \emph{shābāsh}.) A prize? (\textbf{\ur{انعام}}, \emph{inām}.) An award? (\textbf{\ur{اعزاز}}, \emph{ezāz}.) A medal? (\textbf{\ur{تمغہ}}, \emph{tamghā}.) Deserving? (\textbf{\ur{مستحق}}, \emph{mustahiq}.) To tidy up? (\textbf{\ur{سنوارنا}}, \emph{sanvārnā}.) Approval? (\textbf{\ur{منظوری}}, \emph{manzūrī}.) An arrangement? (\textbf{\ur{انتظام}}, \emph{intizām}.) Absolute? (\textbf{\ur{قطعی}}, \emph{qatī}.) To rake in? (\textbf{\ur{بٹورنا}}, \emph{baṭornā}.) To fail an exam? (\textbf{\ur{فیل} \ur{ہونا}}, \emph{fel honā}.) Upbringing? (\textbf{\ur{تربیت}}, \emph{tarbiyat}.) Generous? (\textbf{\ur{فراخ} \ur{دل}}, \emph{farākh-dil}.) Manners? (\textbf{\ur{آداب}}, \emph{ādāb}.) Pampering? (\textbf{\ur{لاڈ}}, \emph{lāḍ}.) To learn by rote? (\textbf{\ur{رٹنا}}, \emph{raṭnā}.) A scholarship? (\textbf{\ur{وظیفہ}}, \emph{vazīfā}.) A difference? (\textbf{\ur{فرق}}, \emph{farq}.) A topic? (\textbf{\ur{موضوع}}, \emph{mauzū}.) Admission? (\textbf{\ur{داخلہ}}, \emph{dākhilā}.) To prefer? (\textbf{\ur{ترجیح} \ur{دینا}}, \emph{tarjīh denā}.)
\end{tcolorbox}
\section[(dialogue)]{Choosing What to Study — once more}

\label{lesson:UR-R385-review-final-385}



\begin{quote}
Say the Urdu for a scholarship and a difference.
\end{quote}

\subsection*{Your turn}
\begin{itemize}
\raggedright
  \item \emph{Say it:} a scholarship, a difference, a topic, admission, to prefer
\end{itemize}

\begin{tcolorbox}[breakable,colback=teal!4,colframe=teal!35!black,title={Before you move on}]
A scholarship? (\textbf{\ur{وظیفہ}}, \emph{vazīfā}.) A difference? (\textbf{\ur{فرق}}, \emph{farq}.) A topic? (\textbf{\ur{موضوع}}, \emph{mauzū}.) Admission? (\textbf{\ur{داخلہ}}, \emph{dākhilā}.) To prefer? (\textbf{\ur{ترجیح} \ur{دینا}}, \emph{tarjīh denā}.)
\end{tcolorbox}
